Richard Pomeroy is a dedicated astrophysics researcher whose academic journey reflects a transition from a successful career in information systems to high-achieving academic scholarship. He graduated Doctor of Philosophy in physics from the University of Texas, Rio Grande Valley (UTRGV, 2026), where he was awarded the President's Research Fellowship. His past research focused on the spatial and statistical characteristics of ultra-compact dwarfs (UCDs) and globular clusters in the Coma Cluster, building upon his long-standing interest in compact stellar systems (CSS).

Richard holds an MSc by Research in Astrophysics from the University of Central Lancashire (UCLan, 2022), where his thesis explored a search for IMBH candidates in CSS using time-domain astronomy and optical emissions. His research incorporated elements such as black hole formation, tidal disruption events, and accretion disc modeling. He previously earned a First-Class Honors BSc in Astronomy from UCLan (2020), where he received the award for Best Performance in Astronomy and completed an undergraduate dissertation on the same theme.

His background in electronic engineering (BEng, Portsmouth Polytechnic, 1987) and an extensive professional career in IT and systems management has equipped him with exceptional technical skills that he now applies to astronomical data analysis and Python-based programming.

Committed to advancing the field of astrophysics, Richard continues to engage with the broader scientific community through membership in the American Astronomical Society (AAS) and the Royal Astronomical Society (FRAS). His journey reflects a deep passion for astronomy and also a unique blend of analytical rigor, cross-disciplinary expertise, and a drive to contribute meaningfully to the understanding of galaxy and cluster formation processes.

Email: \href{mailto:richard@astropom.com}{richard@astropom.com}